% style.tex -- shared preamble for the standalone test reports.
%
% Each report is an independent document: build it with
%   latexmk -xelatex -interaction=nonstopmode -halt-on-error -outdir=build <file>.tex
% from the reports/ directory. The main project report is intentionally not
% part of this directory; these reports document the test evidence on its own.

\documentclass[UTF8,a4paper,11pt]{ctexart}

\usepackage{geometry}
\geometry{top=2.3cm,bottom=2.3cm,left=2.5cm,right=2.5cm,headheight=14pt}

\usepackage{fontspec}
\setmainfont{DejaVu Serif}[AutoFakeSlant=0.2]
\setsansfont{DejaVu Sans}[AutoFakeSlant=0.2]
\setmonofont{DejaVu Sans Mono}[Scale=0.88]

\usepackage{amsmath,amssymb}
\usepackage{graphicx}
\graphicspath{{figures/}}
\usepackage{booktabs}
\usepackage{tabularx}
\usepackage{longtable}
\usepackage{array}
\usepackage{enumitem}
\usepackage{caption}
\usepackage{float}
\usepackage{xcolor}
\usepackage{listings}
\usepackage{fancyhdr}
\usepackage{hyperref}

\definecolor{codebg}{gray}{0.96}
\definecolor{codeframe}{gray}{0.45}
\lstset{
  basicstyle=\ttfamily\small,
  breaklines=true,
  breakatwhitespace=true,
  columns=fullflexible,
  keepspaces=true,
  frame=single,
  framerule=0.3pt,
  rulecolor=\color{codeframe},
  backgroundcolor=\color{codebg},
  showstringspaces=false,
  tabsize=2,
  aboveskip=0.7em,
  belowskip=0.7em,
}

% Inline code. detokenize preserves spaces and special characters; a plain
% lstinline argument loses spaces under ctex (xeCJK-listings turns them into
% CJK glue). Avoid a literal percent sign in the argument.
\newcommand{\code}[1]{\texttt{\detokenize{#1}}}

\pagestyle{fancy}
\fancyhf{}
\fancyhead[L]{\small 迷你搜索引擎项目}
\fancyhead[R]{\small 测试报告}
\fancyfoot[C]{\thepage}
\renewcommand{\headrulewidth}{0.4pt}
\fancypagestyle{plain}{%
  \fancyhf{}%
  \fancyfoot[C]{\thepage}%
  \renewcommand{\headrulewidth}{0pt}%
}

\captionsetup{font=small,labelfont=bf,skip=5pt}
\setlength{\parskip}{0.35em}
\setlength{\parindent}{2em}
\setlength{\emergencystretch}{3em}
\renewcommand{\arraystretch}{1.15}
\setlist{itemsep=0.2em,topsep=0.3em}

\hypersetup{
  colorlinks=true,
  linkcolor=blue!55!black,
  urlcolor=blue!55!black,
  bookmarksnumbered=true,
}

% Report title block: number/title, subtitle, version.
\newcommand{\reportheader}[3]{%
  \begin{center}
    {\large 迷你搜索引擎项目}\\[0.35em]
    {\LARGE\bfseries #1}\\[0.45em]
    {\large #2}\\[0.6em]
    {\small 版本 #3 \quad 测试日期 2026-10-19}
  \end{center}
  \vspace{0.3em}
  \noindent\rule{\linewidth}{0.6pt}
  \vspace{0.8em}
}

% A highlighted conclusion.
\newcommand{\finding}[1]{%
  \begin{quote}
    \noindent\textbf{结论：}#1
  \end{quote}
}


\begin{document}

\reportheader{测试报告 02：正确性与独立对拍}{黄金样例、Python 参考实现与全量 brute-force}{1.0}
\hypersetup{pdftitle={测试报告 02：正确性与独立对拍}}

\begin{center}\small
\begin{tabularx}{0.92\linewidth}{@{}l X@{}}
  报告编号 & TR-02 \\
  被测对象 & 迷你搜索引擎（C++17，\code{mse}） \\
  被测版本 & 分支 \code{test/reports}，基于 \code{origin/main} 加注释与测试改动 \\
  测试环境 & WSL2（Linux 6.18），g++ 15.2，CMake 3.31.6，Python 3 \\
  测试日期 & 2026-10-19 \\
  总体结论 & 黄金样例、词频 CSV、全词表 brute-force、短语与磁盘词典均一致 \\
\end{tabularx}
\end{center}

\section{目标与范围}

本报告回答一个与单元测试不同的问题：\textbf{C++ 实现是否忠实于规格}。单元测试只能证明“代码符合我理解的规格”；如果规格本身被误解，测试也会一起错。因此这里使用三类独立证据：

\begin{enumerate}
  \item \textbf{黄金样例}：讲义给出的四文档例子，人工逐项核对 posting 表；
  \item \textbf{独立参考实现}：Python 重新实现分词、停用词、Porter 词干化与索引解析，从原始文本重建期望值；
  \item \textbf{全量 brute-force}：不只抽查几个词，而是对全部 14,937 个词干逐个比较 df、tf、docID 与 positions。
\end{enumerate}

阈值实验与规模测试分别见 TR-03、TR-04；本报告不重复。

\section{方法}

\subsection{黄金样例}

\code{tests/fixtures/golden4/} 复刻讲义 1.2 节的四篇文档：

\begin{table}[H]
  \centering
  \caption{黄金样例文档}
  \label{tab:golden-docs}
  \begin{tabularx}{\linewidth}{@{}r X@{}}
    \toprule
    docID & 文本 \\
    \midrule
    1 & Gold silver truck \\
    2 & Shipment of gold damaged in a fire \\
    3 & Delivery of silver arrived in a silver truck \\
    4 & Shipment of gold arrived in a truck \\
    \bottomrule
  \end{tabularx}
\end{table}

原始构建（无停用词、无词干化）得到 4 个文档、11 个词、25 个 token。程序输出的 posting 表与讲义逐项一致：

\begin{table}[H]
  \centering
  \caption{黄金样例 posting 表}
  \label{tab:golden-postings}
  \begin{tabularx}{0.72\linewidth}{@{}l X@{}}
    \toprule
    词 & posting（docID, tf） \\
    \midrule
    a & (2,1) (3,1) (4,1) \\
    arrived & (3,1) (4,1) \\
    damaged & (2,1) \\
    delivery & (3,1) \\
    fire & (2,1) \\
    gold & (1,1) (2,1) (4,1) \\
    in & (2,1) (3,1) (4,1) \\
    of & (2,1) (3,1) (4,1) \\
    shipment & (2,1) (4,1) \\
    silver & (1,1) (3,2) \\
    truck & (1,1) (3,1) (4,1) \\
    \bottomrule
  \end{tabularx}
\end{table}

位置校验：doc1 中 \code{silver} 在 position 1；doc3 中 \code{silver} 在 positions 2 和 6。

\textbf{一处参考材料复核}：讲义表格把 \code{a} 记为 (2,1)(4,1)，但 doc3 的原文是 “Delivery of silver arrived in \textbf{a} silver truck”，按句子重算应为 (2,1)(3,1)(4,1)。程序输出与原文一致。测试因此以原始句子为准，并在用例中显式记录这处差异。这说明黄金样例本身也需要复核，不能盲目相信二手表格。

停用词版本（过滤 \code{a/of/in}）得到 8 个词；\code{silver} 在 doc3 中的位置重新编号为 1 和 3，验证了“停用词不占 position”的语义。

\subsection{独立 Python 实现}

三个脚本各自独立实现规格，不调用 C++ 的中间结果：

\begin{itemize}
  \item \code{tools/check_wordcount.py}：独立分词、停用词过滤与 Porter 词干化，生成与 \code{mse count} 同格式的 CSV；
  \item \code{tools/brute_force.py}：从 \code{data/text/} 重新扫描语料，重建完整倒排索引；同时用独立二进制解析器读取 \code{index.bin}（v1 或 v2），逐词比较；
  \item \code{tools/check_phrases.py}：从语料随机抽取短语，在 Python 的规范化 token 流上做暴力相邻匹配，再与 \code{mse query --phrase} 的文档集比较。
\end{itemize}

“独立”指实现路径独立：Python 侧不链接 C++ 代码，也不读取 C++ 的中间结构；只有原始 manifest 与文本层是共同输入。

\section{对拍结果总表}

{\small
\begin{longtable}{@{}>{\raggedright\arraybackslash}p{0.38\linewidth}
                    >{\raggedright\arraybackslash}p{0.24\linewidth}
                    >{\raggedright\arraybackslash}p{0.32\linewidth}@{}}
  \caption{独立对拍结果}\label{tab:crosscheck-all} \\
  \toprule
  对拍项 & 规模 & 结果 \\
  \midrule
  \endfirsthead
  \toprule
  对拍项 & 规模 & 结果 \\
  \midrule
  \endhead
  \bottomrule
  \endlastfoot
  逐文档 token 数（C++ vs Python） & 195 文档 & 100\% 一致，866,626 token \\
  \code{wordcount.csv} 默认配置 & 25,705 行 & 逐字节一致 \\
  \code{wordcount.csv --keep-speakers} & 25,705 行 & 逐字节一致 \\
  \code{wordcount.csv --stop-words} & 556,866 kept / 309,760 removed & 逐字节一致 \\
  \code{wordcount.csv --stem} & 14,947 词干 & 逐字节一致（NLTK \code{MARTIN_EXTENSIONS}） \\
  \code{wordcount.csv --stem --stop-words} & 14,937 词干 & 逐字节一致 \\
  词典随机工作负载 & 20,000 次操作 / 300 词 & splay 与 hash 词表完全一致 \\
  索引文件（splay vs hash） & 3,685,655 B & 逐字节一致 \\
  索引文件（splay vs \code{std::map}） & 3,685,655 B & 逐字节一致 \\
  \code{brute_force.py} 全词表（v1） & 14,937 词干 & df、tf、docID、positions 全项一致 \\
  \code{brute_force.py} 全词表（v2） & 14,937 词干 & 独立解析 v2，全项一致 \\
  \code{brute_force.py} raw 版 & 25,705 词 & 全项一致 \\
  块式构建 vs 内存构建 & 17 个 run / 4,166,556 B 临时 & 索引逐字节一致 \\
  磁盘 B+Tree \code{dump --verify-all} & 14,937 词 & 0 差异 \\
  磁盘词典查询 vs 全量加载 & \code{osric} 单词查询 & 文本与 JSON 输出逐字节一致 \\
  短语查询 vs 独立扫描（v1） & 303 条随机短语 & 返回文档集 100\% 一致 \\
  短语查询 vs 独立扫描（v2） & 303 条随机短语 & 返回文档集 100\% 一致 \\
  RFC 语料全链路 & 3 文档 / 2,660 词干 & count CSV 逐字节一致；brute-force 全项一致 \\
\end{longtable}
}

\section{关键对拍细节}

\subsection{词频 CSV：逐字节一致}

\code{mse count} 与 \code{check_wordcount.py} 在五种配置下生成的文件用 \code{diff} 比较，均无差异。词频 CSV 同时包含 rank、tf、df、$df\_ratio$、$tf\_share$ 与熵特征，因此这项比较不仅验证“词和次数”，也验证了熵公式、排序与浮点格式。

\subsection{全词表 brute-force}

\code{brute_force.py} 不抽样：它对索引中的每一个 term 重建 posting，再逐项比较。正式配置下的输出为：

\begin{lstlisting}[language=,basicstyle=\ttfamily\small]
index    : build/index.bin
documents: 195   terms: 14937   tokens: 556866
compared : 14937 terms (df, tf, docIDs and positions for every posting)
result   : OK -- independent scan and index file agree
\end{lstlisting}

这也覆盖了停用词位置语义：停用词不占 position，因此短语相邻性在过滤后的 token 流上仍然成立。

\subsection{短语对拍}

\code{check_phrases.py} 从语料随机抽取短语（默认 300 条，实际采样 305 条）：

\begin{lstlisting}[language=,basicstyle=\ttfamily\small]
corpus   : data/manifest.tsv (195 documents)
index    : build/index.bin
phrases  : 305 sampled, 303 checked, 2 skipped
result   : OK -- every checked phrase matches the brute-force scan
\end{lstlisting}

两条被跳过：一条在规范化后只剩一个可检索词（退化为单词查询），另一条在索引中没有任何可检索词（全部是停用词）。跳过原因会打印，不是静默忽略。同一脚本对 v1 与 v2 各跑一遍，结果一致。

\subsection{三种词典产出同一索引}

splay、hash、\code{std::map} 三种词典分别构建正式索引，三个文件逐字节相同（3,685,655 B）。随机工作负载测试（20,000 次操作、300 个词）也比较了 splay 与 hash 的词表，结果完全一致。字典实现只影响构建期性能，不影响冻结后的索引内容。

\subsection{v2 压缩索引的独立解析}

v2 把 postings 改为 docID gap、tf、position delta，并用字节对齐 varint 编码。\code{brute_force.py} 内置独立 v2 解析器，不调用 C++ 解码函数；14,937 个词全部一致。另有 varint 边界测试覆盖 6/14/22/30 位的上下界与截断拒绝。

\subsection{磁盘 B+Tree 词典}

正式索引的 sidecar 为 237,568 B（58 页 $\times$ 4 KiB，树高 2）。\code{dump --verify-all} 逐个 term 经 B+Tree 定位并解析 posting，结果：

\begin{lstlisting}[language=,basicstyle=\ttfamily\small]
verified   : 14937 terms, 0 mismatches
dict reads : 114 page reads, 29819 cache hits
\end{lstlisting}

同一索引的全量加载路径与磁盘路径对 \code{osric} 的文本输出、JSON 输出逐字节一致。

\subsection{第二语料：RFC}

\code{data2/} 收录 RFC 1035/5321/7231 三篇 plain 文本。用同一套代码构建（动态熵停用词 58 词 + 词干化）得到 2,660 个词干、4,594 条 posting、45,370 个 token、308,164 B。\code{brute_force.py} 对 2,660 个词全项一致；\code{mse count} 的 CSV 与 Python 版本逐字节一致。这证明引擎不依赖莎士比亚语料的具体格式。

\section{对拍发现并修复的问题}

对拍与测试共抓出 13 个真实问题。下面保留完整清单，因为它们比“全部通过”更能说明测试的有效性。

{\footnotesize
\begin{longtable}{@{}r >{\raggedright\arraybackslash}p{0.19\linewidth}
                    >{\raggedright\arraybackslash}p{0.34\linewidth}
                    >{\raggedright\arraybackslash}p{0.33\linewidth}@{}}
  \caption{测试与对拍发现的问题}\label{tab:all-bugs} \\
  \toprule
  \# & 现象 & 原因 & 处理 \\
  \midrule
  \endfirsthead
  \toprule
  \# & 现象 & 原因 & 处理 \\
  \midrule
  \endhead
  \bottomrule
  \endlastfoot
  1 & 对拍脚本比 C++ 少算 4,340 个 token & 只按行首判断标记行，把 \code{[To HELENA]} 之类台词误判为标题行 & 改为行首与行尾同时匹配 \\
  2 & 正文混入 \code{aside}、\code{sings} 和人名 & 610 处方括号舞台提示被当作正文 & 提取阶段把行首 \code{[...]} 重定位为 \code{(...)} 舞台行；净减 849 token \\
  3 & 过滤版词频的 \code{tf_share} 分母错误 & C++ 曾把过滤前 token 总数记入 \code{tokens} & 只累计通过过滤器的 token；\code{stopwordTokens} 单独累计 \\
  4 & 41 个词形被过度词干化（如 \code{agreem}） & 误以为后缀条件失败后应继续尝试更短后缀 & 改为“首个匹配说了算、条件失败即停” \\
  5 & 伸展树插入后段错误 & 插入分支保留了 \code{Node*&} 子指针引用，旋转后引用失效 & 先把新节点指针复制到局部变量，再 splay \\
  6 & 索引加载误报 “index file is truncated” & \code{Reader} 构造时传入了临时 vector，整句结束后引用悬垂 & 先把文件内容存入具名变量，再构造 Reader \\
  7 & v2 压缩索引的 posting offset 全部偏移 & 计算压缩偏移时漏加了 term 表自身大小 & 复用与 v1 相同的 \code{computeTermOffsets()} \\
  8 & 块式构建在真实语料上报 “run 文件截断” & Windows 上 run 文件用文本模式读取，CRLF 被翻译进二进制流 & 显式 \code{std::ios::binary} \\
  9 & \code{MapDictionary} 的 \code{sortedEntries()} 产出空词表 & \code{try_emplace} 只填 map 的键，\code{TermEntry::term} 未镜像 & 插入后手动设置 \code{entry.term} \\
  10 & B+Tree 的 I2 区间校验误报 & 分隔键是真实键的副本，右子树最左键等于左分隔键 & I2 改为左闭区间 \code{[key[i-1], key[i])} \\
  11 & 变长键下 B+Tree 的 50\% 填充不变量被违反 & 按条数对半分裂在键长偏斜时不保证字节下界 & 改为“合法窗口内取字节中点”分裂 \\
  12 & 磁盘词典查询头部显示 “1 terms” & lazy 路径只把被查询词挂到内存索引 & 头部改用 sidecar 的 \code{keyCount()} \\
  13 & 30M/50M 键的 B+Tree bulk-load 报 I2 违反 & 内层借用时插入了被移动子树的最小键，而非原 \code{last} 的首键 & 插入 \code{last} 原来的首键作为边界 \\
\end{longtable}
}

\section{复现方式}

\begin{lstlisting}[language=bash]
# 词频 CSV 对拍
python3 tools/check_wordcount.py > build/wordcount_py.csv
diff build/wordcount.csv build/wordcount_py.csv

# 全词表 brute-force（v1）
python3 tools/brute_force.py --stop-words data/stopwords.txt --stem

# 短语对拍（v1；把 index 换成 v2 索引即可对拍压缩格式）
python3 tools/check_phrases.py --binary ./build/mse \
    --index build/index.bin --stop-words data/stopwords.txt --stem --queries 300

# 磁盘词典全词表校验
./build/mse dump --index build/index.bin --dict-file build/index.bpt --verify-all

# 需要的索引
./build/mse build --index build/index.bin --dict-file build/index.bpt \
    --stop-words data/stopwords.txt --stem
\end{lstlisting}

上述命令都已接入 CTest（\code{crosscheck_*}），日常回归只需 \code{ctest --test-dir build -R crosscheck}。

\section{结论}

\begin{itemize}
  \item 黄金样例的 posting 表与原始文本一致，并纠正了参考表格中 \code{a} 的一处错误；
  \item 五种词频配置的 CSV 与独立 Python 实现逐字节一致；
  \item 14,937 个词干的 df、tf、docID、positions 全量对拍通过（v1 与 v2）；
  \item 303 条随机短语在 v1/v2 上均与独立扫描一致；
  \item 三种词典实现产出逐字节相同的索引；
  \item 磁盘 B+Tree 的 14,937 个词 0 差异；
  \item RFC 第二语料端到端一致，证明引擎不依赖莎士比亚格式。
\end{itemize}

\finding{在当前语料与当前规格下，C++ 索引与独立 Python 参考实现达到逐字节/逐项一致；测试不是“再跑一遍 C++”，而是从原始文本重建期望值。}

\end{document}
